\sectioncentered*{Заключение}
\addcontentsline{toc}{section}{Заключение}

В ходе работы над данной курсовой работой было произведено существенное расширение  функционала разрабатываемого приложения с добавлением возможности использования AI-чата для быстрых вопросов по истории болезней пациента и рекомендациям по здоровью без привлечения помощи медицинского персонала, а также страницы личных сообщений для для обеспечения оперативной коммуникации между медицинским персоналом и пациентом в режиме реального времени, интегрирована RAG система для загрузки и обработки медицинских документов, а также возможности составления отчетов по состоянию здоровью пациента, добавлена страница мониторинга состояния здоровья для регулярного обновления и актуализации данных о состоянии медицинских показателей здоровья пациента, возможности их дальнейшего прогнозирования, а также их первичной статистической обработки и предоставления лечащему врачу в удобном визуальном формате, кроме того была  существенно улучшена визуальная составляющая и удобство использования приложения для обеспечения интуитивно понятного и функционального доступа пользователей к взаимодействию с разрабатываемой интеллектуальной системой. Был реализован единый, удобный и функциональный интерфейс в виде веб-приложения, представляющий собой минимальный жизнеспособный продукт. Был получен опыт в сфере командной работы над проектом, включающий опыт использования выбранного оптимальным подхода к управлению - 'Метод критического пути', опыт совместного планирования, разработки и реализации проектов в качестве руководителя проекта.

Таким образом, в результате проектирования интеллектуальной медицинской системы была реализована архитектура системы, учитывающая выявленные в ходе анализа требования к функциональности, безопасности, а также масштабируемости системы. В связи с этим данная система позволяет обеспечить удобство и надежность использования как для медицинских специалистов, так и для пациентов.
